\documentclass[10pt,a4paper,journal]{article}
\usepackage[utf8]{inputenc}
\usepackage[brazil]{babel}
\usepackage[margin=2.2cm]{geometry}
\usepackage{amsmath,amsfonts,amssymb}
\usepackage{booktabs}
\usepackage{microtype}
\usepackage{hyperref}
\usepackage{caption}
\usepackage{setspace}
\onehalfspacing

\title{Resumo Expandido Artigo Consad}
\author{Eduardo Lopes Jonker, Vivian Kremer Kalempa, Yuri Kaszubowski Lopes\\ \small Universidade do Estado de Santa Catarina (UDESC) -- PPGCAP}
\date{2026}

\begin{document}
\maketitle

\subsection{Relato de Experiência e Modelagem Baseada em Censo de Frotas do Governo do Estado de Santa Catarina}
**Evento/Destino:** Congresso CONSAD de Gestão Pública (Conselho Nacional de Secretários de Estado da Administração)

**Autor:** Eduardo Lopes Jonker

**Orientadora:** Profa. Dra. Vivian Kremer Kalempa

**Coorientador:** Prof. Dr. Yuri Kaszubowski Lopes

**Instituição:** Universidade do Estado de Santa Catarina (UDESC) - Centro de Educação Superior do Alto Vale do Itajaí (CEAVI) / Centro de Ciências Tecnológicas (CCT)

\begin{abstract}
Este resumo expandido sintetiza o relato de experiência para o CONSAD sobre o uso de IoT e modelos prescritivos para mitigar erros humanos, combater fraudes em combustíveis e otimizar o ciclo de vida de frotas públicas.

**Palavras-chave:** Gestão Pública Prescritiva; Internet das Coisas (IoT); Otimização de Frotas; Auditoria Governamental; Governança Pública.

---

\end{abstract}

\section{1. Problemática}
A administração de frotas no setor público estadual enfrenta desafios crônicos de custos imprevistos em oficinas, vulnerabilidade a desvios no abastecimento de combustíveis e retenção antieconômica de veículos sucateados que oneram o orçamento e reduzem a disponibilidade de viaturas da segurança pública e saúde.

\section{2. Problema Científico}
A transição de uma gestão pública meramente reativa para uma postura prescritiva exige superar a barreira da assimetria de informações e do erro humano. O problema científico consiste em como integrar sensores telemáticos de IoT e ciência de dados a modelos formais de decisão alocativa, fornecendo aos secretários de administração decisões matematicamente blindadas perante os órgãos de controle externo.

\section{3. Objetivo}
O objetivo é relatar a implementação de um modelo prescritivo baseado em IoT e otimização matemática para o controle de manutenção e abastecimento da frota de Santa Catarina, automatizando a glosa de inconsistências e identificando veículos candidatos a alienação sob o Decreto nº 47.539/2018.

\section{4. Método}
Apoiado em microdados censitários de 350.737 itens em 77.867 ordens de serviço (R\$ 138,76 milhões pagos e R\$ 60,05 milhões em glosas) e 1.000.000 de abastecimentos, o método articula: (i) telemetria IoT com leitura em tempo real de barramento CAN e odometria; (ii) algoritmos de detecção de anomalias (Isolation Forest) para identificar desvios de consumo e superfaturamento; (iii) modelagem prescritiva MILP para alocação orçamentária mensal; e (iv) aplicação do critério regulatório dos 40\% da FIPE.

\section{5. Resultado}
O framework proporcionou redução direta de 17,51\% nas despesas de manutenção (economia de R\$ 9.302.147,26 no cenário S0), bloqueio preventivo de fraudes de abastecimento por divergência entre odômetro e volume bombeado, identificação técnica de 313 veículos onerosos para leilão imediato e projeção de economia de R\$ 40,19 milhões no triênio 2027 a 2029, elevando a prontidão operacional da frota para 95,26\%.

\section{Referências Bibliográficas}
\begin{itemize}
\item Barbado, A., \& Corcho, O. (2022). Interpretable anomaly detection in vehicle fleets. Expert Systems with Applications, 201, 117188.
\end{itemize}
\begin{itemize}
\item Bertsimas, D., \& Kallus, N. (2020). From predictive to prescriptive analytics. Management Science, 66(3), 1025-1044.
\end{itemize}
\begin{itemize}
\item Brasil. (2021). Lei Federal nº 14.133, de 1º de abril de 2021. Lei de Licitações e Contratos Administrativos.
\end{itemize}
\begin{itemize}
\item Santa Catarina. (2018). Decreto Estadual nº 47.539, de 03 de abril de 2018.
\end{itemize}
\end{document}